\documentclass[11pt]{article}
\usepackage[margin=1in]{geometry}
\usepackage{times}
\usepackage{hyperref}
\hypersetup{colorlinks=true, linkcolor=black, urlcolor=blue, citecolor=black}
\title{The Ball That Fell Faster: Bruce DePalma's Spin Experiments and the Timing Problem}
\author{Flyxion \\ \small Independent Researcher}
\date{}
\begin{document}
\maketitle

\begin{abstract}
Bruce DePalma's reported 1970s experiments, in which a rapidly spinning steel ball launched vertically was claimed to rise higher and fall in less time than an identical non-spinning ball launched with the same initial velocity, have circulated as evidence that rotation can alter an object's inertial or gravitational mass. This paper argues that the reported result is best explained by systematic timing and launch-mechanism asymmetries between the spinning and non-spinning trials rather than by any genuine alteration of gravitational behavior, that later, better-controlled replications have failed to reproduce the effect, and that the proposed mechanism, insofar as DePalma specified one, does not identify any established channel by which mechanical spin angular momentum could couple to gravitational mass.
\end{abstract}

\section{Introduction}

Bruce DePalma, a former MIT researcher, reported in the mid-1970s that a steel ball spun to high angular velocity and launched vertically from a spring-loaded mechanism rose noticeably higher, and returned to the launch point in noticeably less time, than an otherwise identical non-spinning ball launched from the same mechanism at nominally the same speed. DePalma interpreted this as evidence that rapid spin reduces an object's effective gravitational mass, a claim he connected to a broader proposed theory of rotational physics distinct from standard mechanics.

\section{The Launch Mechanism Problem}

The reported experiment compared a spinning ball to a non-spinning ball launched from the same spring mechanism, but a ball's interaction with a mechanical launcher is not spin-independent: a rapidly spinning ball in contact with a launch cup or track experiences frictional and gyroscopic interactions during the launch phase itself that a non-spinning ball does not, potentially imparting a different effective launch velocity even when the spring mechanism's stored energy is nominally identical between trials. DePalma's original reports do not describe launch-velocity measurement independent of the timing results themselves, meaning the comparison assumes, rather than independently confirms, that both balls left the launcher at the same speed, which is precisely the assumption a spin-dependent launch interaction would violate.

\section{What Controlled Replication Has Found}

Subsequent attempts to replicate the experiment with independently measured launch velocities, including work using high-speed photography to directly measure the balls' velocity immediately after leaving the launcher rather than inferring it from the launch mechanism's nominal specification, have found that spinning and non-spinning balls launched at genuinely identical measured velocities rise to the same height and return in the same time, within measurement error, consistent with ordinary mechanics and inconsistent with DePalma's spin-dependent gravity claim. The original result is thus best explained by an uncontrolled difference in actual launch velocity between the two conditions, masked by the assumption that a nominally identical spring launch produces identical velocity regardless of the projectile's spin state.

\section{Why the Proposed Mechanism Was Never Specified Precisely Enough to Test Independently}

DePalma's broader theoretical writing proposed that rotation constitutes a distinct and previously unrecognized form of energy and inertia not captured by standard mechanics, but did not derive a quantitative prediction for how much gravitational mass reduction a given angular velocity should produce, a prediction that would be needed to distinguish the claim from a null result on statistical grounds rather than relying on a single, methodologically compromised timing comparison. Without a quantitative prediction, the claim cannot be said to have survived the null replications so much as to have never been placed in a position to be tested against them precisely.

\section{The Entropic Gravity Framing}

On an account where gravitational curvature tracks the redistribution of an ordered configuration at scales where mass-energy content is significant, the angular momentum of a spinning steel ball a few centimeters across represents a change in the ball's internal organization utterly negligible relative to what would be required to alter its participation in that redistribution, consistent with, though not dependent on, the general-relativistic expectation that inertial and gravitational mass are equivalent regardless of an object's internal rotational state, a form of the equivalence principle tested to high precision by entirely different experiments with no observed spin dependence.

\section{Conclusion}

The DePalma spinning ball result is a case study in how an uncontrolled variable in a launch mechanism can produce a striking apparent effect that vanishes under independently measured, better-controlled conditions, and the absence of a quantitative theoretical prediction from DePalma's own proposed mechanism means the claim was never precise enough to be meaningfully confirmed or refuted by a single demonstration in the first place.

\end{document}
